\documentclass[aspectratio=169,10pt]{beamer}

% Shared Beamer setup for Fall 2026 course slides.
\mode<presentation>{
  \usetheme{Madrid}
  \usecolortheme{default}
}

\definecolor{CalPolyGreen}{HTML}{154734}
\definecolor{CalPolyGold}{HTML}{BD8B13}
\definecolor{DeckInk}{HTML}{1F2933}
\definecolor{DeckMuted}{HTML}{5F6B7A}

\setbeamercolor{structure}{fg=CalPolyGreen}
\setbeamercolor{palette primary}{bg=CalPolyGreen,fg=white}
\setbeamercolor{palette secondary}{bg=DeckInk,fg=white}
\setbeamercolor{palette tertiary}{bg=CalPolyGold,fg=DeckInk}
\setbeamercolor{frametitle}{bg=CalPolyGreen,fg=white}
\setbeamercolor{title}{fg=CalPolyGreen}
\setbeamercolor{normal text}{fg=DeckInk,bg=white}
\setbeamercolor{block title}{bg=CalPolyGreen,fg=white}
\setbeamercolor{block body}{bg=black!3,fg=DeckInk}

\setbeamertemplate{navigation symbols}{}
\setbeamertemplate{footline}[frame number]
\setbeamertemplate{itemize item}{\color{CalPolyGold}$\blacktriangleright$}
\setbeamertemplate{itemize subitem}{\color{CalPolyGreen}$\bullet$}

\newcommand{\CourseNumber}{Course}
\newcommand{\CourseTitle}{Course Title}
\newcommand{\LectureNumber}{0}
\newcommand{\LectureTitle}{Lecture Title}
\newcommand{\LectureDate}{Date}
\newcommand{\CourseUnit}{Unit}

\newcommand{\coursenumber}[1]{\renewcommand{\CourseNumber}{#1}}
\newcommand{\coursetitle}[1]{\renewcommand{\CourseTitle}{#1}}
\newcommand{\lecturenumber}[1]{\renewcommand{\LectureNumber}{#1}}
\newcommand{\lecturetitle}[1]{\renewcommand{\LectureTitle}{#1}}
\newcommand{\lecturedate}[1]{\renewcommand{\LectureDate}{#1}}
\newcommand{\courseunit}[1]{\renewcommand{\CourseUnit}{#1}}

\author{Kirk Duran}
\institute{California Polytechnic State University, San Luis Obispo}
\date{\LectureDate}

\newcommand{\makedecktitle}{
  \begin{frame}[plain]
    \vspace{0.25in}
    {\usebeamerfont{title}\color{CalPolyGreen}\LectureTitle\par}
    \vspace{0.18in}
    {\large \CourseNumber{}: \CourseTitle\par}
    \vspace{0.08in}
    {\large Lecture \LectureNumber{} -- \LectureDate\par}
    \vfill
    {\small Kirk Duran \textbar{} Cal Poly, San Luis Obispo\par}
  \end{frame}
}

\newcommand{\placeholdernote}{
  \begin{block}{Placeholder}
    Replace these bullets with the final lecture material, examples, and in-class activities.
  \end{block}
}

\AtBeginSection[]{
  \begin{frame}{Roadmap}
    \tableofcontents[currentsection]
  \end{frame}
}


\usepackage{tikz}

\graphicspath{{../assets/lecture-08/}}

% If an image file is not yet available, skip the visual slot.
\newcommand{\lectureimage}[2][1.75in]{%
  \IfFileExists{../assets/lecture-08/#2}{%
    \includegraphics[width=\linewidth,height=#1,keepaspectratio]{#2}%
  }{%
  }%
}

% Temporary visible placeholder used while lecture assets are being assembled.
% Replace each placeholder with \lectureimage[...]{} once the corresponding asset exists.
\newcommand{\lectureplaceholder}[2][1.75in]{%
  \begin{center}
    \fbox{%
      \parbox[c][#1][c]{0.90\linewidth}{%
        \centering\small\textcolor{DeckMuted}{\textbf{Image placeholder}\\[0.08in]#2}%
      }%
    }
  \end{center}%
}

\setbeamersize{text margin left=0.42in,text margin right=0.42in}

\coursenumber{CS 4880}
\coursetitle{Artificial Intelligence}
\lecturenumber{8}
\lecturetitle{Constraint Satisfaction I}
\lecturedate{September 11, 2026}
\courseunit{Local and Constraint Search}

\begin{document}

\makedecktitle

\begin{frame}[shrink=20]{Today}
  \begin{itemize}
    \item Why many problems are better described by \emph{requirements} than by action sequences.
    \item How a constraint satisfaction problem is defined using variables, domains, and constraints.
    \item How unary, binary, higher-order, and global constraints capture problem structure.
    \item How partial assignments and consistency define a specialized search space.
    \item How recursive backtracking systematically constructs a satisfying assignment.
    \item How map coloring and Sudoku become CSPs under the same general formulation.
  \end{itemize}

  \begin{keyidea}
    Constraint satisfaction changes the search abstraction: instead of asking which path reaches a goal, we ask which assignment makes all required relationships true.
  \end{keyidea}
\end{frame}

\sectionframe{Part 1: From General Search to Constraint Satisfaction}

\begin{frame}[shrink=20]{Humans Reason With Constraints Constantly}
  \begin{columns}[T,onlytextwidth]
    \begin{column}{0.55\textwidth}
      Many everyday planning tasks are naturally expressed as restrictions on a final configuration.

      \begin{itemize}
        \item Schedule meetings so that required participants do not overlap.
        \item Seat guests so that incompatible people are separated.
        \item Assign classrooms while respecting capacity and availability.
        \item Color neighboring regions differently.
      \end{itemize}

      \begin{block}{Common structure}
        We know what relationships must hold, but we may not know the sequence of choices that will produce a valid configuration.
      \end{block}
    \end{column}

    \begin{column}{0.39\textwidth}
      \lectureimage[2.55in]{slide4.png}
    \end{column}
  \end{columns}

  \begin{keyidea}
    Constraint problems are defined more naturally by what must be true than by how the solution is reached.
  \end{keyidea}
\end{frame}

\begin{frame}[shrink=20]{Map Coloring as a Constraint Problem}
  \begin{columns}[T,onlytextwidth]
    \begin{column}{0.51\textwidth}
      Suppose every region of a map must be assigned one of three colors.

      \begin{itemize}
        \item Each region receives exactly one color.
        \item Adjacent regions may not receive the same color.
        \item We care about the final legal coloring, not the order in which regions were colored.
      \end{itemize}

      \begin{block}{Question}
        Can Australia be colored using only \{Red, Green, Blue\} while respecting every adjacency?
      \end{block}
    \end{column}

    \begin{column}{0.43\textwidth}
      \lectureimage[2.70in]{slide5.png}
    \end{column}
  \end{columns}
\end{frame}

\begin{frame}[shrink=20]{Scaling Exposes the Real Problem}
  \begin{columns}[T,onlytextwidth]
    \begin{column}{0.54\textwidth}
      Small instances may be solvable by inspection. Larger instances expose the combinatorial structure.

      \begin{itemize}
        \item With $n$ variables and $d$ possible values per variable, there are up to $d^n$ complete assignments.
        \item Most assignments may violate constraints.
        \item Treating every assignment as an opaque state ignores useful relationships among variables.
      \end{itemize}

      \begin{block}{Example}
        Forty-eight states with three candidate colors already permit
        \[
          3^{48}
        \]
        possible complete colorings before constraints are used.
      \end{block}
    \end{column}

    \begin{column}{0.40\textwidth}
      \lectureimage[2.65in]{slide6.png}
    \end{column}
  \end{columns}

  \begin{keyidea}
    The central challenge is not generating assignments; it is exploiting constraints to avoid generating assignments that cannot possibly work.
  \end{keyidea}
\end{frame}

\begin{frame}[shrink=20]{Why General Search Is an Awkward Fit}
  \begin{columns}[T,onlytextwidth]
    \begin{column}{0.47\textwidth}
      \begin{block}{Classical search}
        \begin{itemize}
          \item Must encode partial assignments as states.
          \item Generates many illegal or redundant branches.
          \item Different assignment orders can represent the same partial solution.
          \item Generic successor generation does not automatically exploit constraint structure.
        \end{itemize}
      \end{block}
    \end{column}

    \begin{column}{0.47\textwidth}
      \begin{block}{Local search}
        \begin{itemize}
          \item Can begin from a complete but inconsistent assignment.
          \item May become trapped or fail to establish that no solution exists.
          \item Requires an objective that quantifies violation quality.
          \item Does not naturally provide a complete systematic solver.
        \end{itemize}
      \end{block}
    \end{column}
  \end{columns}

  \begin{keyidea}
    CSP algorithms specialize search around the internal structure of assignments and constraints rather than treating states as indivisible objects.
  \end{keyidea}
\end{frame}

\begin{frame}[shrink=20]{The Declarative Perspective}
  \begin{columns}[T,onlytextwidth]
    \begin{column}{0.55\textwidth}
      A constraint model specifies the conditions that any acceptable solution must satisfy.

      \begin{itemize}
        \item The model describes \emph{what} a solution looks like.
        \item The solver determines \emph{how} to construct one.
        \item The same solving machinery can apply to maps, schedules, puzzles, layouts, and allocation problems.
      \end{itemize}

      \begin{block}{Shift in representation}
        Search is no longer primarily over paths. It is over assignments to named variables with explicit domains and relationships.
      \end{block}
    \end{column}

    \begin{column}{0.39\textwidth}
      \lectureimage[2.50in]{slide8.png}
    \end{column}
  \end{columns}

  \begin{keyidea}
    Specify what must be true, then let the search procedure determine a consistent assignment.
  \end{keyidea}
\end{frame}

\sectionframe{Part 2: The CSP Model --- Variables, Domains, and Constraints}

\begin{frame}[shrink=20]{Constraint Satisfaction Problem}
  \begin{cpdefinition}
    A \textbf{constraint satisfaction problem (CSP)} is a triple
    \[
      (X,D,C),
    \]
    where $X$ is a set of variables, $D$ assigns a domain of possible values to each variable, and $C$ is a set of constraints over those variables.
  \end{cpdefinition}

  \begin{columns}[T,onlytextwidth]
    \begin{column}{0.30\textwidth}
      \begin{block}{Variables}
        \[
          X=\{X_1,\ldots,X_n\}
        \]
        Quantities whose values must be chosen.
      \end{block}
    \end{column}

    \begin{column}{0.30\textwidth}
      \begin{block}{Domains}
        \[
          D_i=\mathrm{Dom}(X_i)
        \]
        Values currently allowed for $X_i$.
      \end{block}
    \end{column}

    \begin{column}{0.34\textwidth}
      \begin{block}{Constraints}
        Relations that specify which combinations of values are permitted.
      \end{block}
    \end{column}
  \end{columns}
\end{frame}

\begin{frame}[shrink=20]{Assignments: The Vocabulary of CSP Search}
  \begin{columns}[T,onlytextwidth]
    \begin{column}{0.47\textwidth}
      \begin{block}{Partial assignment}
        Assigns values to only some variables.
        \[
          \{WA=R,\;NT=G\}
        \]
      \end{block}

      \begin{block}{Complete assignment}
        Assigns a value to every variable in $X$.
      \end{block}
    \end{column}

    \begin{column}{0.47\textwidth}
      \begin{block}{Consistent assignment}
        Violates no constraint whose variables have all been assigned.
      \end{block}

      \begin{block}{Solution}
        A complete assignment that satisfies every constraint in $C$.
      \end{block}
    \end{column}
  \end{columns}

  \begin{keyidea}
    Backtracking search grows a partial consistent assignment until it either becomes a complete solution or reaches a contradiction.
  \end{keyidea}
\end{frame}

\begin{frame}[shrink=20]{Australia as a CSP}
  \begin{columns}[T,onlytextwidth]
    \begin{column}{0.55\textwidth}
      \begin{block}{Variables}
        \[
          X=\{WA,NT,SA,Q,NSW,V,T\}.
        \]
      \end{block}

      \begin{block}{Domains}
        For every region,
        \[
          D_i=\{R,G,B\}.
        \]
      \end{block}

      \begin{block}{Binary constraints}
        Examples include
        \[
          WA\neq NT,\quad WA\neq SA,\quad NT\neq SA,
        \]
        \[
          NT\neq Q,\quad SA\neq Q,\quad Q\neq NSW.
        \]
      \end{block}
    \end{column}

    \begin{column}{0.39\textwidth}
      \lectureimage[2.75in]{slide12.png}
    \end{column}
  \end{columns}
\end{frame}

\begin{frame}[shrink=20]{Constraint Graphs}
  \begin{cpdefinition}
    For a binary CSP, a \textbf{constraint graph} has one vertex for each variable and an edge between two variables when a binary constraint relates them.
  \end{cpdefinition}

  \begin{columns}[T,onlytextwidth]
    \begin{column}{0.50\textwidth}
      \begin{itemize}
        \item Vertices expose the variables that remain to be assigned.
        \item Edges expose direct dependencies among assignments.
        \item Graph structure will later guide variable-ordering and inference algorithms.
      \end{itemize}

      \begin{block}{Australia}
        Tasmania is isolated in the standard adjacency model, while South Australia is highly connected.
      \end{block}
    \end{column}

    \begin{column}{0.44\textwidth}
      \lectureimage[2.60in]{slide13.png}
    \end{column}
  \end{columns}

  \begin{keyidea}
    The graph tells the solver which decisions can directly constrain one another.
  \end{keyidea}
\end{frame}

\begin{frame}[shrink=20]{A Solution Is Not a Path}
  \begin{columns}[T,onlytextwidth]
    \begin{column}{0.50\textwidth}
      One possible solution is a set of variable-value bindings such as
      \[
        WA=R,\quad NT=G,\quad SA=B,
      \]
      \[
        Q=R,\quad NSW=G,\quad V=R,\quad T=R.
      \]

      What matters is that:
      \begin{itemize}
        \item every variable has a value from its domain;
        \item every constraint is satisfied.
      \end{itemize}
    \end{column}

    \begin{column}{0.44\textwidth}
      \lectureimage[2.55in]{slide14.png}
    \end{column}
  \end{columns}

  \begin{keyidea}
    Two algorithms that reach the same complete satisfying assignment have solved the same CSP even if their internal search histories differ.
  \end{keyidea}
\end{frame}

\sectionframe{Part 3: Constraints Have Structure}

\begin{frame}[shrink=20]{Unary and Binary Constraints}
  \begin{columns}[T,onlytextwidth]
    \begin{column}{0.47\textwidth}
      \begin{block}{Unary constraint}
        Restricts one variable.
        \[
          X\neq Red
        \]
        It can often be enforced simply by removing values from the variable's domain.
      \end{block}
    \end{column}

    \begin{column}{0.47\textwidth}
      \begin{block}{Binary constraint}
        Restricts a pair of variables.
        \[
          X\neq Y
        \]
        Binary constraints are naturally represented as edges in a constraint graph.
      \end{block}
    \end{column}
  \end{columns}

  \vspace{0.10in}
  \lectureimage[1.55in]{slide16.png}
\end{frame}

\begin{frame}[shrink=20]{Higher-Order and Global Constraints}
  \begin{columns}[T,onlytextwidth]
    \begin{column}{0.48\textwidth}
      \begin{block}{Higher-order constraint}
        Relates three or more variables simultaneously.

        Example: three assigned courses cannot occupy the same room-time combination.
      \end{block}

      \begin{block}{Global constraint}
        A single high-level relation over a collection of variables.
      \end{block}
    \end{column}

    \begin{column}{0.46\textwidth}
      \begin{block}{Example: \texttt{AllDifferent}}
        \[
          \mathrm{AllDifferent}(X_1,\ldots,X_k)
        \]
        requires every variable in the scope to take a distinct value.
      \end{block}

      This is a natural representation for rows, columns, and boxes in Sudoku.
    \end{column}
  \end{columns}

  \begin{keyidea}
    Constraints are not limited to pairwise inequality; expressive global relations can summarize substantial problem structure.
  \end{keyidea}
\end{frame}

\begin{frame}[shrink=20]{Constraints Beyond ``Not Equal''}
  \begin{columns}[T,onlytextwidth]
    \begin{column}{0.30\textwidth}
      \begin{block}{Precedence}
        One activity must occur before another.
        \[
          start(A)+dur(A)\le start(B)
        \]
      \end{block}
    \end{column}

    \begin{column}{0.31\textwidth}
      \begin{block}{Disjunctive}
        At least one alternative relationship must hold.
        \[
          A\prec B\;\lor\;B\prec A
        \]
      \end{block}
    \end{column}

    \begin{column}{0.33\textwidth}
      \begin{block}{Resource / capacity}
        Multiple decisions share a limited resource such as rooms, machines, or staff.
      \end{block}
    \end{column}
  \end{columns}

  \vspace{0.12in}
  \lectureimage[1.45in]{slide18.png}
\end{frame}

\begin{frame}[shrink=20]{Why Structure Helps Search}
  \begin{columns}[T,onlytextwidth]
    \begin{column}{0.54\textwidth}
      In generic search, a state is often treated as an indivisible object.

      In a CSP, the solver can inspect the internal structure:
      \begin{itemize}
        \item which variables are already assigned;
        \item which values remain legal;
        \item which variables constrain one another;
        \item which assignments immediately violate a relation.
      \end{itemize}

      \begin{block}{Consequence}
        Search can reject an entire family of complete assignments after detecting a contradiction in only a small partial assignment.
      \end{block}
    \end{column}

    \begin{column}{0.40\textwidth}
      \lectureimage[2.55in]{slide19.png}
    \end{column}
  \end{columns}
\end{frame}

\sectionframe{Part 4: Consistency --- Detecting Impossible Values Early}

\begin{frame}[shrink=20]{Local Consistency: The Central Question}
  \begin{columns}[T,onlytextwidth]
    \begin{column}{0.55\textwidth}
      Constraints let us ask whether values remain compatible with what is currently known.

      \begin{block}{Reasoning question}
        Can this candidate value still participate in some assignment that satisfies the relevant constraints?
      \end{block}

      \begin{itemize}
        \item If not, the value can be removed from consideration.
        \item If a variable has no values left, the current search branch is impossible.
        \item Stronger forms of consistency inspect larger neighborhoods of the constraint graph.
      \end{itemize}
    \end{column}

    \begin{column}{0.39\textwidth}
      \lectureimage[2.55in]{slide21.png}
    \end{column}
  \end{columns}

  \begin{keyidea}
    Constraint reasoning converts logical incompatibilities into domain reduction before complete assignments are generated.
  \end{keyidea}
\end{frame}

\begin{frame}[shrink=20]{Node Consistency}
  \begin{cpdefinition}
    A variable is \textbf{node-consistent} when every value remaining in its domain satisfies all unary constraints on that variable.
  \end{cpdefinition}

  \begin{columns}[T,onlytextwidth]
    \begin{column}{0.52\textwidth}
      Suppose
      \[
        D_X=\{1,2,3,4\}
      \]
      with unary constraint
      \[
        X>2.
      \]

      Enforcing node consistency removes values $1$ and $2$:
      \[
        D_X\leftarrow\{3,4\}.
      \]
    \end{column}

    \begin{column}{0.42\textwidth}
      \lectureimage[2.20in]{slide22.png}
    \end{column}
  \end{columns}
\end{frame}

\begin{frame}[shrink=20]{Arc Consistency --- Intuition}
  \begin{cpdefinition}
    For a directed arc $X_i\rightarrow X_j$, the arc is \textbf{consistent} when every value in $D_i$ has at least one supporting value in $D_j$ that satisfies the binary constraint between the variables.
  \end{cpdefinition}

  \begin{itemize}
    \item A value without any compatible neighbor value cannot appear in a solution consistent with that arc.
    \item Removing one value may cause other values elsewhere to lose their support.
    \item This creates the possibility of \emph{constraint propagation}.
  \end{itemize}

  \begin{keyidea}
    Today we use arc consistency only as a reasoning concept; Lecture 9 turns it into an explicit propagation algorithm.
  \end{keyidea}
\end{frame}

\begin{frame}[shrink=20]{Arc Consistency Example}
  \begin{columns}[T,onlytextwidth]
    \begin{column}{0.50\textwidth}
      Let
      \[
        D_X=\{1,2,3\},\qquad D_Y=\{2,3\}
      \]
      with constraint
      \[
        X<Y.
      \]

      For the arc $X\rightarrow Y$:
      \begin{itemize}
        \item $X=1$ has support: $Y=2$ or $3$.
        \item $X=2$ has support: $Y=3$.
        \item $X=3$ has no supporting value.
      \end{itemize}

      Therefore
      \[
        D_X\leftarrow\{1,2\}.
      \]
    \end{column}

    \begin{column}{0.44\textwidth}
      \lectureimage[2.60in]{slide24.png}
    \end{column}
  \end{columns}
\end{frame}

\sectionframe{Part 5: Backtracking Search --- A Complete Baseline Solver}

\begin{frame}[shrink=20]{Backtracking Search}
  \begin{cpdefinition}
    \textbf{Backtracking search} constructs a CSP solution incrementally by assigning one unassigned variable at a time and undoing decisions when the current partial assignment cannot be extended consistently.
  \end{cpdefinition}

  \begin{itemize}
    \item Search state: a partial assignment.
    \item Action: choose an unassigned variable and assign one candidate value.
    \item Rejection rule: do not continue with an assignment that violates a constraint.
    \item Goal: a complete consistent assignment.
  \end{itemize}

  \begin{keyidea}
    Backtracking is depth-first search specialized to the structure of variable assignments.
  \end{keyidea}
\end{frame}

\begin{frame}[shrink=20]{Search Over Partial Assignments}
  \begin{columns}[T,onlytextwidth]
    \begin{column}{0.50\textwidth}
      The root is the empty assignment:
      \[
        \{\}.
      \]

      A path may grow as
      \[
        \{WA=R\}
      \]
      \[
        \{WA=R,NT=G\}
      \]
      \[
        \{WA=R,NT=G,SA=B\}.
      \]

      Any extension that violates a constraint is immediately abandoned.
    \end{column}

    \begin{column}{0.44\textwidth}
      \lectureimage[2.65in]{slide27.png}
    \end{column}
  \end{columns}
\end{frame}

\begin{frame}[shrink=20]{Commutativity Removes Redundant Action Orders}
  \begin{columns}[T,onlytextwidth]
    \begin{column}{0.53\textwidth}
      In ordinary path search, different action orders may describe genuinely different paths.

      In a CSP,
      \[
        WA=R\;\text{ then }NT=G
      \]
      and
      \[
        NT=G\;\text{ then }WA=R
      \]
      produce the same partial assignment.

      \begin{block}{Specialization}
        A backtracking solver chooses exactly one variable at each depth, so it need not explore all permutations of assignment order as separate paths.
      \end{block}
    \end{column}

    \begin{column}{0.41\textwidth}
      \lectureimage[2.55in]{slide28.png}
    \end{column}
  \end{columns}

  \begin{keyidea}
    CSP assignment actions are commutative, allowing a much smaller specialized search tree than a naive action-sequence formulation.
  \end{keyidea}
\end{frame}

\begin{frame}[shrink=20]{Basic Backtracking Pseudocode}
  \begin{block}{Algorithm}
  \ttfamily\small
  BACKTRACK(assignment):\\
  \quad if assignment is complete:\\
  \qquad return assignment\\[0.04in]
  \quad $X \leftarrow$ choose an unassigned variable\\[0.04in]
  \quad for value in DOMAIN($X$):\\
  \qquad if value is consistent with assignment:\\
  \qquad\quad add $X=$value to assignment\\
  \qquad\quad result $\leftarrow$ BACKTRACK(assignment)\\
  \qquad\quad if result $\neq$ failure: return result\\
  \qquad\quad remove $X=$value from assignment\\[0.04in]
  \quad return failure
  \end{block}

  \begin{keyidea}
    This is the correct baseline algorithm. Lecture 9 will change how variables and values are selected and add inference between recursive calls.
  \end{keyidea}
\end{frame}

\begin{frame}[shrink=20]{Worked Example: Beginning the Australia Search}
  \begin{columns}[T,onlytextwidth]
    \begin{column}{0.52\textwidth}
      Assume the fixed variable order
      \[
        WA,NT,SA,Q,NSW,V,T
      \]
      and value order
      \[
        R,G,B.
      \]

      Start:
      \[
        \{\}
      \]

      Try
      \[
        WA=R.
      \]

      Next try
      \[
        NT=R.
      \]

      But $WA\neq NT$, so this branch is rejected immediately.
    \end{column}

    \begin{column}{0.42\textwidth}
      \lectureimage[2.70in]{slide32.png}
    \end{column}
  \end{columns}
\end{frame}

\begin{frame}[shrink=20]{Worked Example: Continuing After Failure}
  \begin{columns}[T,onlytextwidth]
    \begin{column}{0.53\textwidth}
      After rejecting $NT=R$, try
      \[
        NT=G.
      \]

      Then consider $SA$:
      \begin{itemize}
        \item $SA=R$ conflicts with $WA=R$.
        \item $SA=G$ conflicts with $NT=G$.
        \item $SA=B$ is consistent with both.
      \end{itemize}

      The partial assignment becomes
      \[
        \{WA=R,\;NT=G,\;SA=B\}.
      \]

      Search continues recursively from this smaller feasible region.
    \end{column}

    \begin{column}{0.41\textwidth}
      \lectureimage[2.65in]{slide34.png}
    \end{column}
  \end{columns}

  \begin{keyidea}
    A single failed consistency test can discard every complete assignment that extends that inconsistent partial assignment.
  \end{keyidea}
\end{frame}

\begin{frame}[shrink=20]{Completeness Does Not Mean Efficiency}
  \begin{columns}[T,onlytextwidth]
    \begin{column}{0.52\textwidth}
      With finite domains, basic backtracking is complete: if a solution exists, exhaustive search will eventually find one.

      But the worst-case search remains exponential.
      \[
        O(d^n)
      \]
      for $n$ variables with at most $d$ candidate values each.

      \begin{itemize}
        \item Arbitrary variable order can postpone obvious failures.
        \item Arbitrary value order can destroy future flexibility.
        \item Basic consistency checks may detect contradictions much later than necessary.
      \end{itemize}
    \end{column}

    \begin{column}{0.42\textwidth}
      \lectureimage[2.55in]{slide35.png}
    \end{column}
  \end{columns}

  \begin{keyidea}
    Backtracking gives us a complete foundation; the next problem is to make that foundation intelligent enough to scale.
  \end{keyidea}
\end{frame}

\sectionframe{Part 6: Modeling a Larger CSP --- Sudoku}

\begin{frame}[shrink=20]{Sudoku as a Constraint Satisfaction Problem}
  \begin{columns}[T,onlytextwidth]
    \begin{column}{0.50\textwidth}
      A standard $9\times 9$ Sudoku can be modeled with one variable per cell.

      \begin{block}{Variables}
        \[
          X_{r,c},\qquad 1\le r,c\le 9.
        \]
      \end{block}

      \begin{block}{Domains}
        \[
          D_{r,c}\subseteq\{1,2,\ldots,9\}.
        \]
        A given clue fixes the corresponding variable to a singleton domain.
      \end{block}
    \end{column}

    \begin{column}{0.44\textwidth}
      \lectureimage[2.70in]{slide36.png}
    \end{column}
  \end{columns}
\end{frame}

\begin{frame}[shrink=20]{Sudoku Constraints}
  \begin{columns}[T,onlytextwidth]
    \begin{column}{0.52\textwidth}
      Every row, column, and $3\times 3$ box must contain distinct digits.

      \begin{block}{Rows}
        \[
          \mathrm{AllDifferent}(X_{r,1},\ldots,X_{r,9})
        \]
      \end{block}

      \begin{block}{Columns}
        \[
          \mathrm{AllDifferent}(X_{1,c},\ldots,X_{9,c})
        \]
      \end{block}

      The same global constraint is applied to each of the nine boxes.
    \end{column}

    \begin{column}{0.42\textwidth}
      \lectureimage[2.65in]{slide37.png}
    \end{column}
  \end{columns}

  \begin{keyidea}
    Map coloring and Sudoku look very different visually, but both reduce to assignments constrained by explicit relations.
  \end{keyidea}
\end{frame}

\begin{frame}[shrink=20]{Naive Backtracking on Sudoku}
  \begin{columns}[T,onlytextwidth]
    \begin{column}{0.54\textwidth}
      A basic solver can proceed exactly as before:

      \begin{enumerate}
        \item choose an unassigned cell;
        \item try one value from its domain;
        \item reject values that violate a row, column, or box constraint;
        \item recurse;
        \item undo the assignment after failure.
      \end{enumerate}

      \begin{block}{The algorithm is general}
        Nothing fundamental about backtracking changes between map coloring and Sudoku; only the CSP model changes.
      \end{block}
    \end{column}

    \begin{column}{0.40\textwidth}
      \lectureplaceholder[2.55in]{Sudoku search sequence: choose cell, try digit, mark conflict, undo, then try another digit.}
    \end{column}
  \end{columns}
\end{frame}

\begin{frame}[shrink=20]{The Problem With Naive Backtracking}
  \begin{columns}[T,onlytextwidth]
    \begin{column}{0.52\textwidth}
      Suppose three unassigned variables currently have domains
      \[
        D_A=\{2,4,7\},
      \]
      \[
        D_B=\{5\},
      \]
      \[
        D_C=\{1,3,8,9\}.
      \]

      A fixed-order solver may choose $A$ or $C$ before $B$ even though $B$ is almost completely determined.
    \end{column}

    \begin{column}{0.42\textwidth}
      \begin{block}{Two questions for Lecture 9}
        \begin{enumerate}
          \item Which variable should we assign next?
          \item What consequences can we infer immediately after an assignment?
        \end{enumerate}
      \end{block}

      \lectureplaceholder[1.65in]{Three variable boxes with domain sizes 3, 1, and 4; visually emphasize the singleton-domain variable as the obvious choice.}
    \end{column}
  \end{columns}

  \begin{keyidea}
    The baseline solver treats many decisions as arbitrary even when the current domains and constraint graph contain strong information about what should happen next.
  \end{keyidea}
\end{frame}

\begin{frame}[shrink=20]{From Basic Search to Intelligent CSP Solving}
  \begin{center}
  \renewcommand{\arraystretch}{1.35}
  \begin{tabular}{l|l}
    \textbf{Today: Foundation} & \textbf{Next: Efficient CSP Solving} \\\hline
    Variables, domains, constraints & Minimum Remaining Values (MRV) \\
    Constraint graphs & Degree heuristic \\
    Partial and complete assignments & Least Constraining Value (LCV) \\
    Consistency checking & Forward checking \\
    Basic backtracking & Arc consistency and propagation \\
    Map coloring and Sudoku models & Maintaining Arc Consistency (MAC)
  \end{tabular}
  \end{center}

  \begin{keyidea}
    Lecture 8 establishes a complete solver; Lecture 9 improves its search order and adds inference so that contradictions are discovered much earlier.
  \end{keyidea}
\end{frame}

\begin{frame}[shrink=20]{Takeaways}
  \begin{itemize}
    \item A CSP is defined by variables $X$, domains $D$, and constraints $C$.
    \item A solution is a complete assignment that satisfies every constraint.
    \item Constraint graphs expose dependencies among variables in binary CSPs.
    \item Constraint types include unary, binary, higher-order, and global relations such as \texttt{AllDifferent}.
    \item Local consistency asks whether candidate values remain supportable under the known constraints.
    \item Backtracking performs depth-first search over partial assignments and is complete for finite CSPs.
    \item Basic backtracking can still require exponential search because it makes uninformed ordering decisions and performs only limited inference.
  \end{itemize}

  \begin{keyidea}
    The power of the CSP representation is that search can reason about the internal structure of a partial solution; the next lecture turns that structure into aggressive pruning and propagation.
  \end{keyidea}
\end{frame}

\end{document}
